\section{Limitations and alternative explanations}
\label{sec:limitations}

\subsection{The gain is process efficiency, not quality uplift}

The most important limitation is also the simplest. Both real-process transitions tie on champion promotions and retained quality. The child wins because it spends fewer hidden observations. A reader who defines RSI only as recursively increasing task quality should regard the present result as evidence about recursive search-process efficiency, not about that stronger definition.

\subsection{G2 contains a development-adapted stopping threshold}

The selected G2 stops when observed best quality is at least 780. This threshold was produced from the V19 public development corpus. Direct access to world identity was forbidden, and the fresh V19 holdout strongly rejects pure memorization of the 201 development worlds, but the threshold is still development-adapted. V20 then evaluates one real state whose public quality is 953. The result does not show that the same stopping rule transfers to unrelated organism families.

\subsection{One real terminal state}

V20 has two independent physical G1 child evaluations, but the process comparison begins from one real A016 state. This supports a causal state-level claim, not a population estimate. Multiple fresh campaigns are needed before making frequency or robustness claims about the real-process advantage.

\subsection{Repeated use of one evaluator epoch}

The hidden evaluator implementation remains fixed and outside the editable system, but prior aggregate results become public over the course of the project. This can create adaptive pressure against one evaluation authority even when hidden cases are never revealed. A strong L4 experiment should use a new evaluator/task family that was not part of V14--V20 development.

\subsection{External LLM mutation operator}

G2 is not written solely by the prior policy. An external LLM supplies candidate source edits from the frozen bundle. The lineage controls what evidence and search state are exposed and which candidate survives, but the mutation operator itself is externally provided. This is comparable to evolutionary systems that use a foundation model as a mutation engine, but weaker than a fully autonomous closed loop.

\subsection{Manual proposer isolation}

Temporary-chat isolation is procedural. Hashes and transcript schemas prove which bundle and output artifacts entered the laboratory, but they cannot cryptographically prove that the operator selected the specified model/effort or that no outside context entered a chat. Future work should replace this with a provider-attested or locally isolated proposer runner.

\subsection{Project-defined \lthreep}

The term \lthreep is defined by this project. It should not be read as a consensus field-wide RSI level. We use it to prevent a bounded process result from being conflated with stronger concepts such as unbounded recursive intelligence growth.

\subsection{No statistical population inference}

The key gates are deterministic or paired one-shot protocol tests, not samples designed for population-level statistics. The 24-world holdout is useful evidence of generalization over that generator, but it does not justify a claim that G2 will win on essentially all possible search worlds. The paper reports exact observed counts rather than confidence intervals implying a broader sampling frame.
